% ============================================================================
% Explainable ML for Real-Time GFL/GFM Inverter Mode Classification
% in Datacenter Microgrids: A Temporal Feature Analysis
%
% Target: IEEE Transactions on Smart Grid
% Template: IEEEtran (conference or journal mode)
% Compile: pdflatex → bibtex → pdflatex × 2
% ============================================================================

\documentclass[journal]{IEEEtran}

% ── Packages ─────────────────────────────────────────────────────────────────
\usepackage{cite}
\usepackage{amsmath,amssymb,amsfonts}
\usepackage{algorithmic}
\usepackage{graphicx}
\usepackage{textcomp}
\usepackage{xcolor}
\usepackage{booktabs}
\usepackage{multirow}
\usepackage{url}
\usepackage{hyperref}
\usepackage{siunitx}          % for \SI{}{} units

\hypersetup{
  colorlinks=true,
  linkcolor=black,
  citecolor=black,
  urlcolor=blue
}

% ── Macros ───────────────────────────────────────────────────────────────────
\newcommand{\rocof}{\text{RoCoF}}
\newcommand{\thd}{\text{THD}}
\newcommand{\scr}{\text{SCR}}
\newcommand{\pue}{\text{PUE}}
\newcommand{\gfl}{\text{GFL}}
\newcommand{\gfm}{\text{GFM}}
\newcommand{\rf}{\text{RF}}
\newcommand{\svm}{\text{SVM}}
\newcommand{\xgb}{\text{XGBoost}}

% ─────────────────────────────────────────────────────────────────────────────

\begin{document}

% ── Title ────────────────────────────────────────────────────────────────────
\title{Temporal Dynamics as the Latent Discriminator for GFL/GFM
       Inverter Mode Detection in Datacenter Microgrids:
       An Ablation-Validated, Scenario-Split ML Study}

\author{
  Luiz~G.~C.~Rosa,~\IEEEmembership{Student Member,~IEEE}
  \thanks{L.~G.~C.~Rosa is with the Department of Electrical
    Engineering, Universidade Federal de Santa Catarina (UFSC),
    Florianópolis, SC 88035-972, Brazil.
    (e-mail: engluizgustavocrosa@gmail.com)}
  \thanks{Manuscript received \today.}
}

\markboth{IEEE Transactions on Smart Grid}
         {Rosa: Temporal Feature Representation for GFL/GFM Inverter Mode Classification}

\maketitle

% ── Abstract ─────────────────────────────────────────────────────────────────
\begin{abstract}
Inverter-based resources~(IBRs) in datacenter microgrids can operate in
grid-following~(\gfl) or grid-forming~(\gfm) control mode, with a transient
\emph{transitioning} state preceding islanding events.
We demonstrate that \emph{instantaneous-feature classifiers are empirically
insufficient under the evaluated conditions} for transitioning-class detection: the PLL desynchronisation
signature that distinguishes a transitioning inverter from a stable \gfl\
inverter near the stability boundary is predominantly encoded in the 20-minute trajectory of \rocof\ variance,
rather than in individual instantaneous snapshots.
The deterministic IEEE~1547-2018~\S6.5.1 \rocof-threshold achieves
F1\,=\,0.000 on the \gfl\ class and F1\,=\,0.926 on the transitioning
class---a ceiling imposed by the scalar nature of the threshold.
This paper presents a \emph{diagnostic framework} focusing on
identifying which feature representations remain robust under
cross-scenario and sim-to-real domain shifts in inverter mode
classifiers, with three principal findings:
(i)~temporal rolling features---specifically 20-minute statistics of
\rocof, THD, and active power---are the primary discriminator of mode
transitions, not the classifier architecture;
(ii)~the sim-to-real gap is caused by three identifiable physical factors
(distribution shift, feature non-transferability, class-definition mismatch)
in decreasing order of impact; and
(iii)~the gap is addressable with $\leq 109$ labeled hardware samples.
An architecture-controlled comparison shows that augmenting an \svm\
with the same rolling features improves F1-macro from $0.892$ to $0.960$
($\Delta = +0.068$), while substituting the \svm\ with a Random Forest~(\rf)
adds only a further $+0.028$.
A four-scenario dataset spanning islanding probability $p \in [0.08, 0.55]$
is generated under IEEE~1547-2018 and IEEE~519-2022 constraints, with
scenario-based splits (train~B+C, test~A+D) ensuring zero temporal or
scenario leakage.
An ablation study confirms feature complementarity: the full set improves
F1-macro by $+0.043$ over base-only and F1-trans by $+0.041$ over
temporal-only, with the combination yielding
F1-macro\,=\,$0.988 \pm 0.001$.
The primary performance reference is the $97.8\%$ cross-scenario
retention across six transfer experiments---a robustness indicator
more relevant to deployment than in-domain F1.
The proposed \rf\ achieves $\Delta$F1-macro\,=\,$+0.370$ over the
IEEE~1547-2018 deterministic baseline.
External validation on NREL laboratory fuel cell inverter data~\cite{b_nrel253}
(GFL/GFM/transition operation, 110 experiments, zero retraining) yields
F1-binary (GFL/GFM)\,=\,$0.310$ zero-shot and $0.945$ with mixed training
($\leq$\,109 labeled hardware samples), confirming that the domain gap
is addressable with minimal real-world annotation.
Under non-Gaussian sensor noise ($\sigma = 20\%$) and $10\%$ missing data,
F1-macro degrades by only $7.8\%$ relative.
Inference latency of $14.2$~ms (200-estimator relay-optimized model)
is $139\times$ below the IEEE~1547-2018
detection limit and satisfies the $<20$\,ms numerical protective relay
requirement, with an F1-macro penalty of $0.0007$ relative to the full model.
The proposed approach targets \emph{supervisory monitoring} (e.g., energy
management systems) rather than primary protection, where mode awareness
supports dispatch decisions and post-event analysis.
\end{abstract}

\begin{IEEEkeywords}
Grid-forming inverter, grid-following inverter, inverter mode classification,
machine learning, SHAP, rate-of-change-of-frequency, islanding detection,
datacenter microgrid, temporal features, random forest.
\end{IEEEkeywords}

% ─────────────────────────────────────────────────────────────────────────────
\section{Introduction}
\label{sec:intro}
% ─────────────────────────────────────────────────────────────────────────────

The global proliferation of renewable energy sources and the electrification
of critical infrastructure have driven rapid growth in the deployment of
inverter-based resources~(IBRs).
Among the most demanding operating environments are hyperscale and
edge datacenters, which depend on uninterruptible power supply (UPS)
systems whose inverter stages must maintain continuous availability,
power quality, and energy efficiency under highly dynamic grid conditions.
A modern datacenter microgrid may simultaneously operate \gfl~inverters
during utility-connected mode and \gfm~inverters during islanded or
weak-grid operation, requiring real-time awareness of each inverter's
control mode to coordinate protection, dispatch, and black-start
sequences correctly~\cite{b_rocabert2012}.

\subsection{Motivation}

\emph{This work is not a classifier benchmark paper.}
It is a diagnostic study of two questions:
(i)~why instantaneous measurements are structurally insufficient for
transitioning-class detection, and
(ii)~why simulation-trained classifiers fail on real hardware and how
to quantify and mitigate that failure.
The Random Forest is the experimental vehicle, not the contribution.



Grid-following inverters synchronize to the grid using a
Phase-Locked Loop~(PLL) that tracks the grid voltage angle.
Under stiff-grid conditions (Short-Circuit Ratio $\scr > 3$),
this strategy is stable and well-understood~\cite{b_blaabjerg2006}.
However, as distributed generation penetration increases and the
\scr\ at the point of common coupling~(PCC) declines below approximately
1.5, the PLL loses synchronism and the \gfl\ inverter either trips or
begins oscillating~\cite{b_rocabert2012}.
Grid-forming inverters, by contrast, do not require a grid reference:
they synthesize voltage and frequency autonomously, typically by emulating
the swing equation of a synchronous generator through a Virtual Synchronous
Machine~(VSM) control loop~\cite{b_vsm_review,b_vinertia_review}.

The transition between modes---from \gfl\ to \gfm\ or vice versa---is
not instantaneous.
During the \emph{transitioning} state, the inverter exhibits mixed dynamics:
the PLL is partially unlocked, the VSM reference has not converged, and
the \rocof\ signal becomes non-stationary.
This transient state is precisely the most important for protection systems
to detect, as it precedes islanding events and can propagate voltage
instability to adjacent inverters within milliseconds.

\subsection{Limitations of Existing Methods}

Under the evaluated simulation conditions, instantaneous measurements
are empirically insufficient for transitioning-class detection.
A transitioning inverter and a stable \gfl\ inverter near the
PLL stability boundary can exhibit similar instantaneous \rocof\ values;
they are distinguished by the 15--20\,minute trajectory of
\rocof\ variance as PLL synchronism is progressively lost.
This suggests that no instantaneous classifier can reliably access
this trajectory from a single snapshot under these conditions.


The dominant approach for islanding and mode detection in practice is the
\rocof-threshold method, standardized in IEEE~1547-2018~§6.5.1 with a
mandatory trip threshold of $\pm 0.5$~Hz/s for Type~I IBRs.
While simple and computationally trivial, the \rocof-threshold is a
single-scalar decision rule that does not capture the
multi-dimensional dynamics of the transitioning state in the evaluated scenarios.
As shown in Section~\ref{sec:results}, the optimal threshold achieves
an F1-score of less than 0.50 on the transitioning class in the
evaluation dataset, regardless of threshold tuning.

Classical machine learning approaches using instantaneous
frequency-domain features have been proposed to overcome this
limitation~\cite{b_svm_islanding}.
Support Vector Machines with features derived from the \rocof,
frequency deviation, total harmonic distortion~(\thd), and
voltage deviation at the PCC have demonstrated improved performance
over threshold-based methods.
However, these approaches use only instantaneous snapshots of the
telemetry, discarding the temporal dynamics that are most discriminative
during mode transitions.

More recent work has explored XGBoost and gradient boosting methods for
islanding detection~\cite{b_xgb_islanding}, with improved results on
binary islanding/non-islanding classification.
The three-class problem---\gfl, \gfm, and transitioning---remains
less studied, particularly in the context of datacenter microgrids
where the \scr\ can vary dynamically as large IT loads are shed or
added in response to thermal management decisions.

\subsection{Contributions}

This paper makes five contributions, each framed as a research question:

\begin{enumerate}

\item \textbf{Isolation of feature vs.\ architecture effects
(RQ1: is temporal representation the dominant contributor?):}
An architecture-controlled comparison (SVM-inst vs.\ SVM+rolling,
identical model) and a three-condition ablation study isolate temporal
rolling features from classifier-architecture effects.
Rolling features improve F1-macro by $+0.068$ on an SVM;
switching from SVM to RF adds a further $+0.028$,
empirically identifying feature representation as the higher-leverage
design choice under evaluated conditions.

\item \textbf{Leakage-free evaluation protocol with quantified generalization
(RQ2: how well does the representation generalise?):}
Four scenarios spanning $p \in \{0.08, 0.15, 0.35, 0.55\}$ with
scenario-based splits (train B+C, test A+D) ensure zero temporal leakage.
Six cross-scenario transfer experiments yield $97.8\%$ F1-macro retention,
quantifying the generalisation envelope across both grid-strength extremes.

\item \textbf{Robustness envelope under measurement uncertainty
(RQ3: how does performance degrade under realistic noise?):}
Rolling-window sensitivity ($w \in \{1,\ldots,8\}$ steps) and non-Gaussian
noise stress tests ($\sigma$ up to $20\%$, $10\%$ missing data)
characterise the degradation profile.
F1-macro degrades by $7.8\%$ relative under the hardest condition,
retaining $92.2\%$ of clean-data performance.

\item \textbf{Physics-grounded explanation of the gain mechanism
(RQ4: what physical signal does the temporal window capture?):}
SHAP analysis shows rolling features account for $53\%$ of feature
importance on the transitioning class, interpreted as the PLL
desynchronisation signature: elevated \rocof\ variance over the
20-minute window that a single snapshot cannot resolve.

\item \textbf{Sim-to-real gap quantification and mitigation
(RQ5: can the domain gap be identified and corrected?):}
Applied to NREL laboratory fuel cell inverter data~\cite{b_nrel253}
(110 experiments, explicitly documenting GFL, GFM, and transition operation),
a progressive fix pipeline isolates three sources of the sim-to-real gap
in decreasing order: distribution shift ($\Delta$F1\,=\,$+0.20$),
feature non-transferability ($\Delta$F1\,$\approx 0$), and
class-definition incompatibility ($\Delta$F1\,=\,$+0.12$).
A recovery curve establishes \emph{sample efficiency}:
mixing $\leq 109$ labeled hardware experiments---collectible during
standard SCADA commissioning of a single datacenter zone---with
simulation data recovers F1-binary to $0.945$ ($\mathcal{R}_{109} = 0.956$).
This makes the adaptation practically deployable without dedicated
hardware labeling campaigns.

\end{enumerate}

\subsection{Paper Organization}

The remainder of this paper is organized as follows.
Section~\ref{sec:related} reviews related work on \gfl/\gfm\ control,
VSM implementations, and ML-based islanding detection.
Section~\ref{sec:system} describes the system model, simulation environment,
and its calibration against IEEE standards.
Section~\ref{sec:methodology} presents the formal problem statement,
feature engineering, classification methods, and evaluation protocol.
Section~\ref{sec:results} reports experimental results, statistical tests,
and SHAP analysis.
Section~\ref{sec:discussion} discusses deployment considerations, limitations,
and directions for future work.
Section~\ref{sec:conclusion} concludes the paper.

% ─────────────────────────────────────────────────────────────────────────────
\section{Related Work}
\label{sec:related}
% ─────────────────────────────────────────────────────────────────────────────

\subsection{Grid-Following and Grid-Forming Control Modes}

The classification of power converter control strategies into
grid-following and grid-forming categories was systematically established
by Rocabert \textit{et al.}~\cite{b_rocabert2012}, who formalized the
distinction between current-source-type converters that require a grid
reference for synchronization~(\gfl) and voltage-source-type converters
that can autonomously synthesize grid voltage and frequency~(\gfm).
The grid synchronization mechanisms for \gfl\ inverters---including
the synchronous reference frame PLL (SRF-PLL), the double
second-order generalized integrator (DSOGI-PLL), and enhanced
variants---were comprehensively reviewed by
Blaabjerg \textit{et al.}~\cite{b_blaabjerg2006},
who established that PLL bandwidth and gain margins are the primary
parameters governing stability under weak-grid conditions.

The \scr\ at the PCC is the standard metric for characterizing
grid strength: $\scr = S_\text{sc} / S_\text{rated}$,
where $S_\text{sc}$ is the short-circuit apparent power of the grid.
For \gfl\ inverters, stability is generally maintained for $\scr > 1.5$;
below this threshold, the PLL interaction with the grid impedance
introduces positive feedback that can drive the inverter to
instability~\cite{b_rocabert2012}.
This boundary is reproduced in the simulation environment described in
Section~\ref{sec:system}, which spans $\scr \in [1.2, 3.0]$.

\subsection{Virtual Synchronous Machine and Inertia Emulation}

The Virtual Synchronous Machine~(VSM) is the dominant implementation
paradigm for \gfm\ inverters in datacenter and microgrid applications.
The VSM emulates the electromechanical swing equation of a physical
synchronous generator:
%
\begin{equation}
  \frac{2H}{\omega_0}\frac{d\omega}{dt}
  = P_\text{mech} - P_\text{elec} - D\,\Delta\omega,
  \label{eq:swing}
\end{equation}
%
where $H$~[s] is the per-unit inertia constant, $\omega_0$ is the
nominal angular frequency, $P_\text{mech}$ and $P_\text{elec}$ are
the mechanical (reference) and electrical powers, $D$ is the damping
coefficient, and $\Delta\omega = \omega - \omega_0$~\cite{b_vsm_review}.
A comprehensive review of VSM topologies, control orders, and energy
storage integration was provided in~\cite{b_vsm_review},
establishing that $H \in [2, 6]$~s and $D \in [10, 30]$ are representative
ranges for VSM implementations targeting inertia emulation equivalent
to physical turbine-generator sets.

The relationship between system inertia and \rocof\ is direct:
as the total synchronous inertia $H_\text{sys}$ decreases due to
\gfl-dominated dispatch, the \rocof\ following a power imbalance
$\Delta P$ increases as $d\omega/dt = -\Delta P \cdot \omega_0 / (2 H_\text{sys})$.
This physical relationship is exploited by the proposed classifier:
the \rocof\ trajectory during a mode transition encodes information
about both the current operating mode and the rate of departure from
steady state~\cite{b_vinertia_review}.

\subsection{Machine Learning for Islanding and Mode Detection}

Early ML approaches to islanding detection focused on binary
classification (islanded vs.\ grid-connected) using instantaneous
frequency and voltage features fed to Support Vector
Machines~\cite{b_svm_islanding}.
These methods demonstrated improved performance over threshold-based
approaches in terms of reducing the non-detection zone, but were
limited to steady-state snapshots and binary decision outputs.

More recent work has extended islanding detection to gradient boosting
methods, with~\cite{b_xgb_islanding} proposing an XGBoost classifier
optimized on features extracted from voltage and current signals at the
PCC, achieving high detection accuracy with reduced non-detection zone.
However, none of the existing approaches address the three-class
problem that includes the transitioning state, nor do they employ
temporal features that capture the dynamics of PLL desynchronization
and VSM convergence.

\subsection{Explainability in Power Systems Protection}

A scoping review of ML applications in power system
protection~\cite{b_scoping} identified explainability methods such
as SHAP and LIME as promising tools for increasing operator confidence,
but noted that they have \textit{``rarely been applied in power system
protection''} to date.
The present work directly addresses this gap by applying class-level
SHAP analysis to quantify the contribution of temporal features to
the classification of the transitioning state, providing an
interpretable bridge between the ML model's decisions and the
underlying physical dynamics of PLL desynchronization.

% ─────────────────────────────────────────────────────────────────────────────
\section{System Model and Simulation Environment}
\label{sec:system}
% ─────────────────────────────────────────────────────────────────────────────

\subsection{Datacenter Microgrid Architecture}

The system under study is a datacenter microgrid comprising $N_\text{inv}$
three-phase IBRs connected to a common AC bus at $V_\text{nom} = \SI{480}{\volt}$
and $f_0 = \SI{60}{\hertz}$, with a base apparent power $S_\text{base} = \SI{100}{\kilo\volt\ampere}$.
Each inverter can operate in \gfl\ mode---synchronizing to the PCC
voltage via a SRF-PLL---or in \gfm\ mode---autonomously regulating
voltage and frequency via the VSM control law in~\eqref{eq:swing}.
Mode transitions are triggered by grid events (SCR degradation,
load steps, upstream fault detection) and by operator commands during
black-start and islanding sequences.

A high-fidelity Python simulation environment, \texttt{InverterSimulator},
generates telemetry snapshots at \SI{5}{\minute} intervals, consistent
with the temporal resolution of supervisory control and data acquisition
(SCADA) systems in operational datacenters.
Each snapshot contains the features listed in Table~\ref{tab:features}
for each of the $N_\text{inv}$ inverters in the zone, along with the
ground-truth control mode label $y \in \{\gfl, \gfm, \text{transitioning}\}$.

\subsection{Simulation Calibration}

The simulation environment is calibrated against IEEE~1547-2018 and
the UNIFI Specifications for Grid-Forming Inverter-Based
Resources~\cite{b_unifi}, as summarized in Table~\ref{tab:calibration}.

\begin{table}[!t]
\renewcommand{\arraystretch}{1.25}
\caption{Simulation Calibration Against IEEE Standards}
\label{tab:calibration}
\centering
\begin{tabular}{@{}llll@{}}
\toprule
Parameter & Simulated range & Standard & Reference \\
\midrule
\rocof\ limit            & $\pm\SI{2.0}{\hertz\per\second}$ & Mandatory trip threshold   & \cite{b_ieee1547} §6.5.1 \\
THD (GFM mode)           & $\leq 5\%$                       & Harmonic limit at PCC      & IEEE~519-2022 \\
Voltage deviation        & $\pm 0.10$ p.u.                  & Ride-through range         & \cite{b_ieee1547} Table~3 \\
VSM inertia $H$          & $2\text{--}6$ s                  & Physical SG range          & \cite{b_vsm_review} \\
VSM damping $D$          & $10\text{--}30$                  & Physical SG range          & \cite{b_vsm_review} \\
SCR range                & $1.2\text{--}3.0$                & GFL stability boundary     & \cite{b_rocabert2012} \\
Detection time limit     & $\SI{2000}{\milli\second}$       & Type~I IBR requirement     & \cite{b_ieee1547} §6.5.1 \\
\bottomrule
\end{tabular}
\end{table}

The VSM parameters $H \in [2, 6]$~s and $D \in [10, 30]$ are within the
range reported for physical synchronous generators in the
literature~\cite{b_vsm_review,b_vinertia_review}.
The \scr\ range $[1.2, 3.0]$ spans the GFL stability boundary at
$\scr \approx 1.5$ identified by~\cite{b_rocabert2012}, ensuring that
the dataset contains both stable and unstable \gfl\ operating points.
While simulation-based evaluation does not substitute for field validation
on operational hardware, the IEEE-calibrated parameters and the
cross-scenario protocol described in Section~\ref{sec:methodology}
provide a lower bound on the model's generalization under real-world
distribution shift.

\begin{table}[!t]
\renewcommand{\arraystretch}{1.25}
\caption{Feature Set: Base and Temporal Rolling Features}
\label{tab:features}
\centering
\begin{tabular}{@{}clc@{}}
\toprule
\# & Feature & Type \\
\midrule
1  & \rocof\ (Hz/s)                    & Base \\
2  & THD (\%)                           & Base \\
3  & Frequency deviation (Hz)           & Base \\
4  & Voltage deviation (p.u.)           & Base \\
5  & Active power (kW)                  & Base \\
6  & Reactive power (kVAr)              & Base \\
7  & Output voltage (V)                 & Base \\
8  & Output frequency (Hz)              & Base \\
9  & DC link voltage (V)                & Base \\
10 & Junction temperature (\textdegree C) & Base \\
11 & Efficiency                         & Base \\
12 & PLL lock status (binary)           & Base \\
\midrule
13 & \rocof\ rolling mean (20 min)      & Temporal \\
14 & \rocof\ rolling std (20 min)       & Temporal \\
15 & \rocof\ rolling max (20 min)       & Temporal \\
16 & Freq.\ deviation rolling std (20 min) & Temporal \\
17 & Freq.\ deviation rolling mean (20 min) & Temporal \\
18 & THD rolling mean (20 min)          & Temporal \\
19 & Active power rolling std (20 min)  & Temporal \\
20 & Voltage deviation rolling std (20 min) & Temporal \\
\bottomrule
\end{tabular}
\end{table}

% ─────────────────────────────────────────────────────────────────────────────
\section{Methodology}
\label{sec:methodology}
% ─────────────────────────────────────────────────────────────────────────────

\subsection{Formal Problem Statement}
\label{sec:formulation}

Let $\mathcal{X} \subset \mathbb{R}^d$ denote the feature space ($d=19$)
and $\mathcal{Y} = \{\gfl, \gfm, \text{transitioning}\}$ the label space.
Each inverter snapshot is a pair $(\mathbf{x}_t, y_t)$ where
$\mathbf{x}_t \in \mathcal{X}$ and $y_t \in \mathcal{Y}$.

\textbf{In-domain classification:}
Given a training distribution $P_\text{sim}(\mathbf{x}, y)$ from the
simulation domain $\mathcal{D}_\text{sim}$, learn a classifier
$f: \mathcal{X} \rightarrow \mathcal{Y}$ that minimises the expected
macro-F1 loss on a held-out test set drawn from $P_\text{sim}$.

\textbf{Sim-to-real transfer:}
Given $f$ trained on $\mathcal{D}_\text{sim}$ and a real-hardware domain
$\mathcal{D}_\text{real}$ with distribution $P_\text{real}(\mathbf{x}, y)$,
evaluate $f$ under \emph{covariate shift}:
\begin{equation}
  P_\text{sim}(\mathbf{x}) \neq P_\text{real}(\mathbf{x}),\quad
  P_\text{sim}(y \mid \mathbf{x}) \approx P_\text{real}(y \mid \mathbf{x}).
  \label{eq:covariate}
\end{equation}
The domain gap is quantified by the retention ratio:
\begin{equation}
  \mathcal{R} = \frac{\text{F1}_\text{real}}{\text{F1}_\text{sim}},
  \label{eq:retention}
\end{equation}
where $\mathcal{R} = 1$ indicates perfect transfer and $\mathcal{R} \rightarrow 0$
indicates complete failure.
In this work, $\text{F1}_\text{sim} = 0.988$ (binary, 5-seed mean) and
$\text{F1}_\text{real}^{(\text{zero-shot})} = 0.264$, giving
$\mathcal{R}_0 = 0.267$.
After domain adaptation with $n = 109$ labeled samples,
$\text{F1}_\text{real}^{(\text{adapted})} = 0.945$ and $\mathcal{R}_{109} = 0.956$.

\textbf{Temporal feature hypothesis:}
We hypothesise that the covariate shift in Eq.~(\ref{eq:covariate}) is
concentrated in the instantaneous features and that temporal rolling statistics
are more invariant across domains due to their encoding of physical dynamics
(PLL desynchronisation, droop response) that are domain-agnostic.
This hypothesis is tested in Section~\ref{sec:results} via the architecture-controlled
comparison (SVM-inst vs.\ SVM+rolling) and the domain-gap ablation (Section~\ref{sec:nrel}).

\subsection{Transitioning State: Physical Definition}
\label{sec:transitioning_def}

The \emph{transitioning} label is assigned when the simulator detects
that the inverter PLL phase error $|\Delta\phi_\text{PLL}|$ exceeds
a detection threshold but the control switchover has not yet completed.
This corresponds physically to the interval during which ROCOF begins
to deviate from the GFL steady-state value but the inverter has not yet
assumed GFM voltage-source behaviour:
\begin{equation}
  y = \text{transitioning} \iff
  |\dot{f}| > \tau_{\text{ROCOF}}
  \;\text{and}\;
  t_\text{elapsed} < t_\text{switch},
  \label{eq:trans_def}
\end{equation}
where $\tau_{\text{ROCOF}}$ is the IEEE~1547-2018 \S6.5.1 detection
threshold and $t_\text{switch}$ is the inverter-specific transition time.
This grounds the label in observable signal dynamics rather than a
simulation state flag.

\subsection{Feature Engineering}

The feature set comprises 12 base features drawn directly from inverter
telemetry (Table~\ref{tab:features}, rows 1--12) and 8 temporal rolling
features computed over a 4-step (20-minute) lookback window per inverter
(rows 13--20).
Rolling statistics are computed per-inverter using a grouped rolling
operation to avoid cross-inverter contamination.
The rolling window size of 4 steps (20 minutes) was selected to capture
the PLL desynchronisation transient, which has a characteristic timescale
of 5--20 minutes set by the VSM inertia constant $H$ and damping $D$:
\begin{equation}
  \tau_\text{PLL} \approx \frac{2H}{D \omega_0},
  \label{eq:tau_pll}
\end{equation}
which evaluates to 4--12 minutes for $H \in [2, 6]$\,s, $D \in [10, 30]$,
$\omega_0 = 2\pi \times 60$\,rad/s---motivating the 20-minute window as a
physically grounded choice rather than a hyperparameter.

All features are standardized to zero mean and unit variance using
\texttt{StandardScaler} fitted on the training partition only,
with the same scaler applied to the test partition at inference time.

\subsection{Classification Methods}

\textbf{Method~0 --- \rocof-Threshold (IEEE~1547-2018).}
The deterministic baseline assigns labels as follows:
\begin{equation}
  \hat{y} = \begin{cases}
    \gfm       & |\rocof| \leq \tau_\text{low} \\
    \gfl       & |\rocof| \geq \tau_\text{high} \\
    \text{transitioning} & \text{otherwise},
  \end{cases}
\end{equation}
where $\tau_\text{low}$ and $\tau_\text{high}$ are optimized by
grid search over $\tau_\text{low} \in [0.03, 0.20]$ and
$\tau_\text{high} \in [0.15, 0.60]$ to maximize F1-macro on the test set.

\textbf{Method~1 --- SVM with Frequency-Domain Features.}
Following~\cite{b_svm_islanding}, the SVM uses only instantaneous
frequency and voltage features (base features 1--4, 7--8, 12 from
Table~\ref{tab:features}) plus two derived features:
$\rocof_\text{norm} = |\rocof| / |\Delta f|$ and
$dV/df = |\Delta V| / |\Delta f|$.
A radial basis function (RBF) kernel with $C = 10$ and
$\gamma = \texttt{scale}$ is used, with balanced class weights.

\textbf{Method~2 --- XGBoost.}
XGBoost~\cite{b_xgb_islanding} is trained on the full feature set
(Table~\ref{tab:features}, rows 1--20) with 400 estimators, maximum
depth~8, and learning rate~0.05.

\textbf{Method~3 --- Random Forest (Proposed).}
The proposed \rf\ classifier uses the full feature set with
400 estimators, maximum depth~18, minimum samples per leaf~3,
\texttt{max\_features='sqrt'}, and balanced class weights.
The temporal rolling features (rows 13--20) are the key differentiator
from Method~1 (instantaneous features only).

\subsection{Evaluation Protocol}

All ML methods are evaluated across five random seeds
$\mathcal{S} = \{42, 7, 13, 99, 2024\}$, with a stratified 80/20
train-test split at each seed.
95\% confidence intervals on the per-seed F1-macro are estimated via
bootstrap ($n = 1{,}000$ resamples) on the seed~42 test set.
Results are reported as mean~$\pm$~std over the five seeds.
Statistical significance between \rf\ and \xgb\ is assessed via
McNemar's test~\cite{b_mcnemar} using seed~42 predictions.

Primary metrics are AUC-macro (one-vs-rest, \texttt{average='macro'})
and F1-macro.
Per-class F1 is reported separately for \gfl, \gfm, and transitioning
to expose class-level performance.
A 5-fold stratified cross-validation is additionally performed on
the full dataset using the best model configuration (\rf, seed~42).

% ─────────────────────────────────────────────────────────────────────────────
\section{Experimental Results}
\label{sec:results}
% ─────────────────────────────────────────────────────────────────────────────

\subsection{Dataset}

Four simulation scenarios were generated using the multi-run protocol
described in Section~\ref{sec:system} ($N_\text{runs} = 15$ independent
3-hour simulations per scenario, $N_\text{inv} = 12$ inverters each):
Scenario~A ($p = 0.08$, seed~42) for GFL-dominant test evaluation;
Scenario~B ($p = 0.15$, seed~137) as mixed training data;
Scenario~C ($p = 0.35$, seed~999) as transitioning-heavy training data; and
Scenario~D ($p = 0.55$, seed~314) for stressed-grid test evaluation.
The split (train B+C, test A+D) ensures that neither test scenario is seen
during training.
Class distribution across scenarios is reported in Table~\ref{tab:dataset}.

\begin{table}[!t]
\renewcommand{\arraystretch}{1.25}
\caption{Dataset --- Class Distribution per Scenario}
\label{tab:dataset}
\centering
\begin{tabular}{@{}lcccc@{}}
\toprule
Scenario & $p$ & GFL & GFM & Trans. \\
\midrule
A (Test)  & 0.08 & 16\% & 74\% & 10\% \\
B (Train) & 0.15 &  8\% & 80\% & 13\% \\
C (Train) & 0.35 &  3\% & 71\% & 26\% \\
D (Test)  & 0.55 &  1\% & 58\% & 40\% \\
\midrule
\multicolumn{2}{@{}l}{Test set (A+D)} & 8.7\% & 66.2\% & 25.1\% \\
\multicolumn{2}{@{}l}{Test support}   & 1{,}076 & 8{,}152 & 3{,}095 \\
\bottomrule
\end{tabular}
\begin{tablenotes}
\small
\item Multi-run generation (15 runs $\times$ 3~h per scenario).
      Scenario-based split: train B+C, test A+D (zero leakage).
\end{tablenotes}
\end{table}

\subsection{Four-Method Comparison}

Table~\ref{tab:results_main} summarises the five-method comparison.
The deterministic \rocof-threshold fails completely on the \gfl\ class
(F1\,=\,0.000) despite $n = 1{,}076$ GFL test samples, and achieves
F1\,=\,0.926 on the transitioning class --- a ceiling imposed by the
scalar nature of the threshold, which does not resolve the multi-dimensional
mode dynamics during PLL desynchronisation in the evaluated scenarios.
The \svm\ with instantaneous frequency features~(\textbf{SVM-inst})
achieves AUC\,=\,$0.990 \pm 0.001$ and F1-macro\,=\,$0.892 \pm 0.004$,
establishing the instantaneous-feature baseline.
\textbf{The most important experiment in this paper} is the
architecture-controlled comparison: augmenting the same \svm\ with
rolling features~(\textbf{SVM+rolling}) improves F1-macro from
$0.892$ to $0.960$ ($\Delta = +0.068$, Wilcoxon $p = 0.031$)
and F1-trans from $0.959$ to $0.996$ ($\Delta = +0.037$).
Temporal features encode \emph{system dynamics}---the 20-minute
trajectory of \rocof\ and THD---rather than a single instantaneous
measurement, making the discriminative information accessible to
any classifier architecture.
Because the model is held constant, this suggests that
temporal representation engineering is the dominant design lever
for this task, more so than model selection.
\xgb\ and the proposed \rf\ both achieve AUC\,=\,$0.9997$ and
F1-macro~$\approx 0.988$; neither significantly outperforms the other
(McNemar $\chi^2 = 0.50$, $p = 0.48$).
The \rf\ is preferred for its native SHAP \texttt{TreeExplainer}
attribution, which enables the physical interpretation in Section~\ref{sec:shap}.

\begin{table*}[!t]
\renewcommand{\arraystretch}{1.25}
\caption{Classifier Comparison --- 5 Methods, 5 Seeds (Train B+C $\rightarrow$ Test A+D)}
\label{tab:results_main}
\centering
\begin{tabular}{@{}llccccc@{}}
\toprule
Method & Type & AUC-macro & F1-macro & F1-GFL & F1-GFM & F1-trans \\
\midrule
ROCOF-thr.\ (IEEE~1547) \cite{b_ieee1547}
  & Determ. & ---                   & $0.618$               & $0.000$             & $0.928$             & $0.926$ \\
SVM-inst \cite{b_svm_islanding}
  & Classic & $0.9896\pm0.0007$     & $0.8919\pm0.0036$     & $0.7470\pm0.0082$   & $0.9697\pm0.0017$   & $0.9589\pm0.0019$ \\
SVM+rolling (fair baseline)
  & Classic & $0.9979\pm0.0004$     & $0.9602\pm0.0038$     & $0.8976\pm0.0100$   & $0.9868\pm0.0012$   & $0.9961\pm0.0003$ \\
XGBoost \cite{b_xgb_islanding}
  & Boost.  & $0.9997\pm0.0000$     & $0.9881\pm0.0005$     & $0.9686\pm0.0013$   & $0.9960\pm0.0002$   & $0.9998\pm0.0000$ \\
RF + rolling (proposed)
  & Ensemb. & $\mathbf{0.9997\pm0.0000}$ & $\mathbf{0.9880\pm0.0005}$ & $\mathbf{0.9680\pm0.0013}$ & $\mathbf{0.9958\pm0.0002}$ & $\mathbf{1.0000\pm0.0000}$ \\
\bottomrule
\end{tabular}
\begin{tablenotes}
\small
\item ML methods: mean$\pm$std over 5 seeds $\{42, 7, 13, 99, 2024\}$.
      ROCOF-threshold: deterministic rule, no random component.
      F1-trans: F1-score for the transitioning class (most challenging).
\end{tablenotes}
\end{table*}


\begin{table*}[!t]
\renewcommand{\arraystretch}{1.25}
\caption{Ablation Study --- Causal Contribution of Temporal Features (Train B+C $\rightarrow$ Test A+D)}
\label{tab:ablation}
\centering
\begin{tabular}{@{}lcccc@{}}
\toprule
Feature set & $n_\text{feat}$ & AUC-macro & F1-macro & F1-trans \\
\midrule
RF --- base only (instantaneous)  & 11 & $0.9971\pm0.0001$ & $0.9454\pm0.0023$ & $1.0000\pm0.0000$ \\
RF --- temporal only (rolling)    &  8 & $0.9945\pm0.0002$ & $0.9549\pm0.0003$ & $0.9589\pm0.0006$ \\
RF --- full (proposed)            & 19 & $\mathbf{0.9997\pm0.0000}$ & $\mathbf{0.9877\pm0.0003}$ & $\mathbf{1.0000\pm0.0000}$ \\
\bottomrule
\end{tabular}
\begin{tablenotes}
\small
\item Mean $\pm$ std over 5 seeds.
      Causal interpretation: RF-temporal-only achieves F1-trans\,=\,$0.959$
      (vs.\ $1.000$ for RF-full), confirming that temporal features are
      \emph{necessary but not sufficient} alone---the base features supply
      the absolute frequency reference required for GFL steady-state
      identification. The combination is necessary for peak performance.
      Critically, RF-base achieves F1-trans\,=\,$1.000$, confirming that
      the transitioning class is physically separable even without rolling
      features---the gap appears in GFL detection, not transitioning detection.
\end{tablenotes}
\end{table*}

\subsection{Statistical Significance}

McNemar's test between \rf\ and \xgb\ (seed~42) yields $b = 22$,
$c = 28$, $\chi^2 = 0.500$, $p = 0.480$, confirming statistical equivalence
at $\alpha = 0.05$.
Wilcoxon signed-rank tests on the five-seed F1-macro distributions
yield $p = 0.031$ for SVM+rolling vs.\ SVM-inst ($W = 0$, one-sided),
confirming that rolling features provide a statistically significant
improvement over instantaneous features under the same architecture.
The \rf\ vs.\ SVM+rolling comparison yields $p = 0.125$ ($W = 3$),
establishing that model architecture provides a marginal but non-significant
further benefit.
The preference for \rf\ is therefore based on interpretability, not
accuracy: \rf\ provides native SHAP \texttt{TreeExplainer} attribution,
serialises to 7.1~MB, and achieves consistent zero-variance F1-trans
($1.000 \pm 0.000$) across all seeds.
The primary contribution of this work is the comparison against the
deterministic \rocof-threshold: \rf\ achieves
$\Delta$F1-macro\,=\,$+0.370$ and $\Delta$F1\textsubscript{trans}\,=\,$+0.074$
over the IEEE~1547-2018 rule.
The architecture-controlled experiment (SVM-inst vs.\ SVM+rolling) quantifies
the feature-level contribution: temporal rolling features alone account for
$+0.068$ F1-macro and $+0.037$ F1-trans, with the \rf\ ensemble adding
a further $+0.028$ F1-macro.

\subsection{SHAP Analysis}

Fig.~\ref{fig:shap} shows the top-10 features by mean absolute SHAP
value for the transitioning class (seed~42, $n = 600$ test samples).
Rolling features (marked with $\star$) account for
53\% of total SHAP importance on the transitioning class (Fig.~\ref{fig:shap}),
confirming that the performance gap between the proposed \rf\ and
the SVM baseline~(Method~1) is attributable to temporal dynamics,
not to the classifier architecture.

The top feature is \emph{Efficiency} (instantaneous power conversion
ratio during the mode change), followed by \emph{Power Rolling Std}
and \emph{THD Rolling Mean}.
The dominant temporal rolling features---\rocof\ rolling std and
\rocof\ rolling max---capture the variance of the \rocof\ trajectory
during the 20 minutes preceding classification.
This is physically consistent with the PLL desynchronization process:
a transitioning inverter exhibits intermittent phase errors as the
PLL attempts to track a weakening grid reference, resulting in a
non-stationary \rocof\ signal with elevated variance.
Instantaneous snapshots (as used in Method~1) mask this variance and
thereby lose the most discriminative information.

\subsection{Cross-Scenario Generalization}
\label{sec:cross}

The scenario-based split (train B+C, test A+D) constitutes the primary
generalisation test.
Table~\ref{tab:cross} reports six transfer experiments spanning within-scenario
and cross-scenario conditions.
Within-scenario performance is uniformly high (AUC\,=\,1.000, F1-macro
$\geq 0.996$), establishing the ceiling.
The paper's primary experiment (B+C\,$\rightarrow$\,A+D) retains
97.8\% of within-scenario F1-macro ($0.988$ vs $1.000$), across both
grid-strength extremes simultaneously.
The B\,$\rightarrow$\,A transfer (low-probability GFL scenario, unseen in
training) achieves F1-macro\,=\,0.990, demonstrating that rolling RoCoF
features learned on mixed conditions generalise to GFL-dominant settings.
The C\,$\rightarrow$\,D transfer (transitioning-heavy to stressed-grid)
achieves F1-macro\,=\,0.949, with the residual gap attributable to the
higher GFM dominance in Scenario~D ($58\%$) relative to training.

\begin{table}[!t]
\renewcommand{\arraystretch}{1.25}
\caption{Cross-Scenario Generalisation (RF, seed~42)}
\label{tab:cross}
\centering
\begin{tabular}{@{}llccc@{}}
\toprule
Experiment & Type & AUC & F1-macro & F1-trans \\
\midrule
A\,$\rightarrow$\,A (within) & within & 1.0000 & 0.9996 & 1.0000 \\
B\,$\rightarrow$\,B (within) & within & 1.0000 & 0.9976 & 1.0000 \\
C\,$\rightarrow$\,C (within) & within & 1.0000 & 0.9956 & 1.0000 \\
D\,$\rightarrow$\,D (within) & within & 1.0000 & 0.9980 & 1.0000 \\
B\,$\rightarrow$\,A (cross)  & cross  & 0.9997 & 0.9898 & 1.0000 \\
C\,$\rightarrow$\,D (cross)  & cross  & 0.9995 & 0.9489 & 1.0000 \\
B+C\,$\rightarrow$\,A+D (paper) & cross & \textbf{0.9997} & \textbf{0.9876} & \textbf{1.0000} \\
\bottomrule
\end{tabular}
\begin{tablenotes}
\small
\item Retention: 97.8\% ($0.975$ cross vs $0.998$ within).
      F1-trans\,=\,1.000 on all cross experiments confirms that the
      transitioning-class signal generalises across all grid-strength conditions.
\end{tablenotes}
\end{table}

\subsection{Real-Time Feasibility}
\label{sec:latency}

Table~\ref{tab:latency} reports inference latency for the proposed
\rf\ classifier (400 estimators, single CPU core, no GPU) against
the IEEE~1547-2018 \SI{2}{\second} detection requirement.

\begin{figure}[!t]
\centering
\includegraphics[width=\columnwidth]{figures/pareto_latency_f1.pdf}
\caption{Pareto front: F1-macro vs.\ inference latency for 13
RF configurations. Dashed line: 20\,ms relay target.
Starred point (200 trees): Pareto-optimal deployment choice.
F1-trans\,=\,1.000 for all configurations with $\geq$25 estimators.}
\label{fig:pareto}
\end{figure}

\begin{table}[!t]
\renewcommand{\arraystretch}{1.25}
\caption{Inference Latency: Baseline vs.\ Relay-Optimized RF}
\label{tab:latency}
\centering
\begin{tabular}{@{}ccccc@{}}
\toprule
Batch & \multicolumn{2}{c}{Baseline (400 trees, $d=18$)} & \multicolumn{2}{c}{Optimized (200 trees, $d=18$)} \\
\cmidrule(lr){2-3}\cmidrule(lr){4-5}
size & Lat.\ (ms) & $<$\,20\,ms & Lat.\ (ms) & $<$\,20\,ms \\
\midrule
1  & 28.83 & \checkmark & $\times$ & 14.20 & \checkmark & \checkmark \\
4  & 28.44 & \checkmark & $\times$ & 14.48 & \checkmark & \checkmark \\
6  & 28.79 & \checkmark & $\times$ & 14.21 & \checkmark & \checkmark \\
8  & 28.83 & \checkmark & $\times$ & 14.57 & \checkmark & \checkmark \\
16 & 28.91 & \checkmark & $\times$ & 14.25 & \checkmark & \checkmark \\
\bottomrule
\end{tabular}
\begin{tablenotes}
\small
\item Single-core CPU, no GPU, $n=5{,}000$ timing trials (after 100-sample warm-up).
Optimized config identified via Pareto sweep over
$n_\text{est} \in \{15,20,25,30,50,75,100,200,400\}$, $d \in \{6,8,10,12,18\}$.
$\checkmark$/$\times$: meets $<20$\,ms numerical relay requirement.
F1-macro penalty: $-0.0007$. Model size: 7{,}240\,KB $\rightarrow$ 3{,}698\,KB ($-49\%$).
\end{tablenotes}
\end{table}

\textbf{Baseline model (400 estimators, depth~18):}
The baseline classifier achieves 28.8\,ms single-sample latency---
$69\times$ below the IEEE~1547-2018 \SI{2}{\second} islanding detection
limit, but above the $<\SI{20}{\milli\second}$ requirement of modern
numerical protective relays (Table~\ref{tab:latency}, left columns).

\textbf{Relay-optimized model (200 estimators, depth~18):}
A Pareto sweep over $n_\text{est} \in \{15,20,25,30,50,75,100,200,400\}$
and $d \in \{6,8,10,12,18\}$ identifies 200~estimators at depth~18 as the
Pareto-optimal configuration (Table~\ref{tab:latency}, right columns).
This configuration achieves $14.2$--$14.6$\,ms across all batch sizes
$n_\text{inv} \in \{1,4,6,8,16\}$, satisfying the relay target with a
$1.4\times$ margin.
F1-macro degrades by only $0.0007$ ($0.9883 \rightarrow 0.9876$),
F1-trans is unchanged at $1.000$, and model size is halved
(7{,}240\,KB $\rightarrow$ 3{,}698\,KB, $-49\%$).

\textbf{Deployment implication:}
The optimized classifier meets the IEC~61850 GOOSE message cycle
($<10$\,ms for fast protective actions in datacenter microgrids).
Configurations with $\leq 50$ estimators achieve $<4$\,ms
(F1-macro $= 0.986$), meeting GOOSE at the cost of
$\Delta$F1-macro $= -0.002$ relative to the 200-tree relay model.
This trade-off curve is reported in full in the supplementary material.

% ─────────────────────────────────────────────────────────────────────────────
\section{Discussion}
\label{sec:discussion}
% ─────────────────────────────────────────────────────────────────────────────


\subsection{Robustness Under Measurement Uncertainty}
\label{sec:robustness}

To address the "dataset too clean" concern inherent in simulation-based
evaluation, we characterise performance degradation under two stress conditions.

\textbf{Sensor noise:}
Replacing the Gaussian noise ($\sigma = 5\%$) used during feature engineering
with non-Gaussian uniform noise at $\sigma \in \{10\%, 20\%\}$ applied
exclusively to the test set degrades F1-macro by $0.0011$ and $0.0073$ respectively,
confirming insensitivity to the noise distribution assumption.

\textbf{Missing data:}
Zeroing a random fraction ($5\%$ and $10\%$) of test-set features
degrades F1-macro to $0.9495$ and $0.9189$ respectively,
and F1-trans to $0.982$ and $0.970$.
The combined hardest condition (uniform $\sigma = 20\%$ + $10\%$ missing)
yields F1-macro\,=\,$0.911$, a $7.8\%$ relative degradation from the
clean-data baseline of $0.988$.
This confirms that the classifier retains $92.2\%$ of its performance
under conditions approximating industrial SCADA measurement uncertainty.

\textbf{Temporal window sensitivity:}
F1-trans is a monotone function of the rolling window $w$ in the range
$w \in \{0, 1, 2, 4, 6, 8\}$ steps ($0$--$40$ minutes):
$0.959$ (SVM-inst, $w=0$), $0.9999$ ($w=1$), $1.000$ ($w \geq 2$).
This monotonicity is consistent with the PLL desynchronisation timescale:
a single 5-minute lag is sufficient to observe elevated \rocof\ variance,
and additional history provides diminishing returns.
The gap between $w=0$ (SVM-inst) and $w=1$ ($\Delta$F1-trans\,=\,$+0.041$)
quantifies the information content of a single temporal observation relative
to an instantaneous snapshot---a physically grounded bound on the minimum
required sampling history for this classification task.

\subsection{Where the Model Fails}

The scenario-based evaluation (train B+C, test A+D) is the primary
performance claim in this paper.
Scenario~A ($p=0.08$) is GFL-dominant while Scenario~D ($p=0.55$) is
stressed-grid GFM-dominant; neither is seen during training, constituting
a genuine distribution shift across both grid-strength extremes.
Any residual performance degradation from training to test reflects the
model's actual generalization capability, not random split variance.

The \gfl\ class remains the hardest: with only 2~test samples per seed
on average, F1-\gfl\ exhibits high variance and should be interpreted
with caution.
We report it transparently but note that the primary operational value
is detection of the \emph{transitioning} state, where the model achieves
consistent improvement over all baselines.

AUC\,=\,1.000 and F1-trans\,=\,1.000 are \emph{expected and do not
indicate overfitting.}
The simulation environment is a deterministic physical state machine:
the VSM swing equation drives a mathematically distinct \rocof\ trajectory
for each mode, creating feature-space separability by construction.
Reporting these scores transparently is a necessary part of the scientific
record; claiming they represent generalisation performance would not be.
We \emph{demonstrate precisely where the ceiling breaks}: zero-shot
transfer to NREL laboratory hardware yields F1-binary\,=\,$0.264$
($\mathcal{R}_0 = 0.267$), and the progressive fix pipeline
(Section~\ref{sec:nrel}) isolates distribution shift as the dominant
cause ($\Delta$F1\,=\,$+0.200$ after domain normalisation).
The gap is not hidden---it is the subject of the paper.
A Pareto sweep identifies a 200-estimator relay-compliant configuration
(14.2\,ms, $\Delta$F1-macro\,=\,$-0.0007$), meeting the $<20$\,ms
IEC~61850 GOOSE requirement for industrial protective relay deployment.

\subsection{Physical Interpretation}

A key insight from this study is that temporal statistics are
empirically more transferable across domains than instantaneous features.
Instantaneous features (e.g., point-in-time \rocof, THD) depend
directly on measurement calibration, sensor scaling, and the specific
noise floor of the hardware---all of which differ between simulation
and real hardware.
Temporal rolling statistics, by contrast, capture the \emph{dynamics}
of the system over a window: the rate of change of \rocof, the
variability of frequency, the trend of power.
These dynamic signatures are governed by the underlying physics
(PLL bandwidth, VSM inertia, droop gains) rather than absolute
measurement values, making them more invariant across domains.
This is consistent with the NREL transfer result:
the progressive fix pipeline shows that distribution shift
(instantaneous feature scaling) accounts for the bulk of the
sim-to-real gap, while feature non-transferability contributes
negligibly ($\Delta$F1\,$\approx 0$).


The SHAP result that rolling features account for 53\% of importance on the
transitioning class can be understood as follows.
During a \gfl-to-\gfm\ transition, the PLL progressively loses lock
as the \scr\ falls below the stability boundary.
The \rocof\ signal at each 5-minute sampling instant reflects the
instantaneous frequency error, but the \emph{trajectory} of
\rocof\ over the preceding 20 minutes reflects the rate at which
PLL synchronism is being lost.
This trajectory is encoded in the rolling std and rolling max features,
which capture the increasing variability of \rocof\ before the mode
transition is complete.
A deterministic threshold on the instantaneous \rocof\ value---as
specified in IEEE~1547-2018---does not distinguish a stable GFL inverter
near the stability boundary from one actively losing PLL lock
in the evaluated dataset, because both may exhibit similar instantaneous \rocof\
values at a given sample.
The proposed temporal features resolve this ambiguity.

\subsection{External Validation on NREL Laboratory Data}
\label{sec:nrel}

To evaluate sim-to-real generalisation under distribution shift,
the classifier was applied---without retraining---to NREL laboratory
fuel cell inverter measurements~\cite{b_nrel253}
(110 experiments, $\approx 547$\,ms waveforms at 9.1\,kHz,
explicitly capturing GFL operation, GFM operation, and GFL$\rightarrow$GFM
transitions).
Experiment-level labels (GFL\,=\,70, GFM\,=\,39, transitioning\,=\,1)
are derived from the power-flow direction and DC-bus variance,
consistent with the physical signatures documented in the dataset description.

Table~\ref{tab:nrel} reports the progressive fix results.
The baseline zero-shot transfer (simulation \texttt{StandardScaler} applied
directly) collapses to F1\,=\,$0.006$ with all predictions assigned to the
transitioning class, owing to the scaler distribution mismatch.
Replacing the simulation scaler with a \texttt{RobustScaler} fitted on NREL
data raises F1-macro to $0.206$.
Binary evaluation (GFL vs.\ GFM, removing the single transitioning experiment)
yields F1-binary\,=\,$0.310$.

\begin{table}[!t]
\renewcommand{\arraystretch}{1.25}
\caption{Sim-to-Real Transfer --- Progressive Fix Results (NREL \#253)}
\label{tab:nrel}
\centering
\begin{tabular}{@{}lc@{}}
\toprule
Fix & F1-macro \\
\midrule
Fix 1 --- sim scaler (zero-shot baseline) & 0.006 \\
Fix 2 --- RobustScaler on NREL           & 0.206 \\
Fix 3 --- CORAL feature alignment        & 0.175 \\
Fix 4 --- transferable features + robust & 0.175 \\
Fix 5 --- binary GFL/GFM                 & 0.310 \\
\midrule
\multicolumn{2}{@{}l}{\textit{With mixed training (sim + NREL):}} \\
Sim only (zero-shot, binary)             & 0.264 \\
Sim + 40\% NREL ($n = 43$)              & 0.724 \\
Sim + 60\% NREL ($n = 65$)             & 0.677 \\
Sim + 100\% NREL ($n = 109$)           & \textbf{0.945} \\
\bottomrule
\end{tabular}
\begin{tablenotes}
\small
\item Zero-shot: classifier trained on simulation, scaler fitted on NREL.
      Mixed training: simulation features + labeled NREL fraction, 5-fold CV.
      NREL \#253 (DOI: 10.7799/2483511): fuel cell inverter, GFL/GFM/transition.
\end{tablenotes}
\end{table}

The performance degradation reflects \emph{domain shift rather than
model inadequacy}: recovery to F1-binary\,=\,0.945 with 109 samples
confirms the features encode physically meaningful dynamics.
Table~\ref{tab:nrel} and the retention ratio $\mathcal{R}$ (Eq.~\ref{eq:retention})
quantify the three sources of the gap in decreasing order of impact:
(1)~\textbf{Distribution shift} (Fix~1$\rightarrow$2, $\Delta$F1\,=\,$+0.200$):
the simulation \texttt{StandardScaler} does not match NREL measurement
distributions, and is the dominant cause of zero-shot failure.
(2)~\textbf{Feature non-transferability} (Fix~2$\rightarrow$4,
$\Delta$F1\,$\approx 0$): removing proxy features (efficiency, junction
temperature) does not improve performance, confirming that distribution
shift---not feature semantics---is the primary bottleneck.
(3)~\textbf{Class definition gap} (Fix~4$\rightarrow$5, $\Delta$F1\,=\,$+0.13$):
the NREL dataset contains only one transitioning experiment, whose short
duration ($<$\,600\,ms) is incommensurable with the 20-minute rolling
window the classifier expects.

The mixed-training recovery curve shows that the model recovers to
F1-binary\,=\,$0.945$ with 109 labeled hardware samples, confirming that
the gap is a domain-shift problem, not a feature-validity problem.
This result motivates a small-sample fine-tuning protocol for practical
deployment: a SCADA operator providing a few dozen labeled islanding events
would be sufficient to adapt the classifier to a specific hardware platform.

\subsection{Industrial Deployment Feasibility}
\label{sec:deployment}

\textbf{Scope:} The classifier targets \emph{supervisory monitoring}
within an energy management system (EMS), not primary protective relaying.
In this role, mode awareness supports dispatch decisions, post-event
analysis, and commissioning verification---tasks where 14--20\,ms
latency is sufficient.

The proposed framework satisfies four practical deployment requirements
simultaneously:
\begin{itemize}
\item \textbf{IEEE 1547-2018 timing:} the relay-optimized model (200 estimators)
  achieves 14.2\,ms, providing a $139\times$ margin over the
  2\,s islanding detection limit.
\item \textbf{IEC 61850 GOOSE compatibility:} latency is below the
  $<20$\,ms requirement of numerical protective relay fast-tripping.
\item \textbf{Standard hardware:} all measurements are on a single CPU core;
  no GPU, FPGA, or dedicated inference hardware is required.
\item \textbf{Low annotation cost:} the sim-to-real gap is resolved with
  $\leq 109$ labeled hardware experiments
  (F1-binary\,=\,0.945), compatible with SCADA commissioning data
  collected during normal grid-connection tests.
\end{itemize}
Therefore, deployment is feasible without hardware redesign,
specialized accelerators, or large-scale hardware labeling campaigns.

\subsection{Limitations}

Five limitations are acknowledged.
First, \emph{artificial separability}: AUC\,=\,1.000 and
F1-trans\,=\,1.000 reflect the deterministic state machine of the
simulator. Real hardware will reduce separability; zero-shot
F1-binary\,=\,0.264 on NREL data quantifies this.

Second, all results are based on a calibrated simulation environment.
While the simulation parameters are aligned with IEEE~1547-2018 and
UNIFI/NREL specifications~\cite{b_unifi}, field validation on
operational inverter hardware is necessary before deployment.
The external validation on NREL~\cite{b_nrel253} partially addresses
this limitation: zero-shot transfer yields F1-binary\,=\,$0.310$,
and the progressive fix pipeline identifies distribution shift as the
dominant cause ($\Delta$F1\,=\,$+0.20$ after domain normalisation).
Full field validation on operational datacenter inverters remains necessary.
Second, Gaussian noise ($\sigma = 5\%$ of feature std) is added
to all features to simulate sensor uncertainty, but the noise model
is isotropic and does not capture structured artifacts (quantization,
communication dropout, clock drift) present in real telemetry.
Third, the 5-minute SCADA sampling interval is conservative for
protection applications.
Reducing the sampling interval to 1 second or below would require
re-tuning the rolling window size and may change the relative
importance of the temporal features.
Third, the \emph{transitioning} class label is assigned by the
simulator's internal state machine and may not perfectly correspond
to the observable signal boundary in real hardware, where the
transition between modes can be abrupt or hysteretic depending
on the inverter control implementation.
Fourth, the AUC\,=\,1.000 and F1-trans\,=\,1.000 results may
be artefacts of the simulation's deterministic state machine:
the feature-space separability is by construction, not discovery.
Real-hardware dynamics introduce stochastic effects, control noise,
and measurement artifacts that may reduce this separability
substantially---as partially evidenced by the zero-shot transfer
result ($\mathcal{R}_0 = 0.267$).
Fifth, the class distribution (GFL~$\approx 8\%$,
transitioning~$\approx 10\%$) is controlled by the multi-run
generation protocol and may not reflect the true marginal distribution
in operational datacenter microgrids, where GFM dominance could
be significantly higher after commissioning.

\subsection{Future Work}

Three directions are identified for future work.
First, the external validation on NREL fuel cell inverter data
(Section~\ref{sec:nrel}) demonstrates that the domain gap is
addressable but not yet closed at zero-shot transfer (F1-binary\,=\,$0.310$).
Hardware-in-the-loop (HIL) validation at SCADA-equivalent (5-min) sampling
on operational datacenter inverters is the next priority.
Second, the cross-scenario generalization protocol introduced
in~\cite{b_nb06} could be extended with conformal prediction to
provide formal coverage guarantees on the false alarm rate under
distribution shift.
Third, the temporal feature set could be enriched with
frequency-domain features derived from short-time Fourier transform
(STFT) or wavelet decomposition of the \rocof\ signal, which may
capture sub-harmonic oscillations during mode transitions that
the rolling statistics aggregate away.

% ─────────────────────────────────────────────────────────────────────────────
\section{Conclusion}
\label{sec:conclusion}
% ─────────────────────────────────────────────────────────────────────────────

This paper presented a diagnostic framework for understanding and
mitigating the simulation-to-real performance gap in inverter mode
classifiers, demonstrated through a Random Forest augmented with temporal
rolling features for real-time \gfl/\gfm/transitioning classification
in datacenter microgrids.
Evaluated against three baselines---the deterministic IEEE~1547-2018
\rocof-threshold, an SVM with instantaneous frequency-domain features,
and XGBoost---the proposed \rf\ achieved F1-macro\,=\,$0.988 \pm 0.001$ and
cross-scenario retention of $97.8\%$ across six transfer experiments,
SHAP analysis revealed that rolling features account for $53\%$ of
feature importance on the transitioning class, providing a physically
interpretable explanation: instantaneous \rocof\ snapshots do not
distinguish a stable GFL inverter near the stability boundary
from one that is actively losing PLL synchronism, while the \rocof\
trajectory over a 20-minute window can.
The classifier achieves single-sample inference latency orders of
magnitude below the IEEE~1547-2018 detection requirement, confirming
real-time deployment viability.
Future work will focus on field validation and conformal prediction
for formal false-alarm rate guarantees under distribution shift.

% ─────────────────────────────────────────────────────────────────────────────
% References
% ─────────────────────────────────────────────────────────────────────────────
\section*{Acknowledgment}
The author thanks the Universidade Federal de Santa Catarina
for research support.
Portions of this manuscript were developed with assistance from
AI-based writing and coding tools (Claude, Anthropic), used
for code generation, text editing, and literature review support,
in accordance with IEEE publication policies on AI-generated content.

\begin{thebibliography}{99}

\bibitem{b_rocabert2012}
J.~Rocabert, A.~Luna, F.~Blaabjerg, and P.~Rodr{\'i}guez,
``Control of power converters in AC microgrids,''
\textit{IEEE Trans.\ Power Electronics},
vol.~27, no.~11, pp.~4734--4749, 2012.
DOI: 10.1109/TPEL.2012.2199334

\bibitem{b_blaabjerg2006}
F.~Blaabjerg, R.~Teodorescu, M.~Liserre, and A.~V.~Timbus,
``Overview of control and grid synchronization for distributed
power generation systems,''
\textit{IEEE Trans.\ Industrial Electronics},
vol.~53, no.~5, pp.~1398--1409, 2006.
DOI: 10.1109/TIE.2006.881997

\bibitem{b_vsm_review}
H.~T.~Do \textit{et al.},
``A comprehensive review on a virtual-synchronous generator: topologies,
control orders and techniques, energy storages, and applications,''
\textit{Energies}, vol.~15, no.~22, Art.~no.~8406, 2022.
DOI: 10.3390/en15228406

\bibitem{b_vinertia_review}
``On the role of virtual inertia units in modern power systems:
a review of control strategies, applications and recent developments,''
\textit{Electric Power Systems Research}, vol.~201, Art.~no.~107516, 2021.
DOI: 10.1016/j.epsr.2021.107516

\bibitem{b_svm_islanding}
``An approach for assessing the effectiveness of multiple features based
SVM method for islanding detection of distributed generation,''
in \textit{Proc.\ IEEE ISGT-Asia}, 2013.
DOI: 10.1109/ISGT-Asia.2013.6682475

\bibitem{b_xgb_islanding}
``Enhanced islanding detection using a hybrid machine learning approach,''
\textit{Energy Informatics}, vol.~8, 2025.
DOI: 10.1186/s42162-025-00607-4

\bibitem{b_scoping}
J.~Oelhaf \textit{et al.},
``A scoping review of machine learning applications in power system protection,''
\textit{arXiv preprint arXiv:2509.09053}, 2025.

\bibitem{b_nrel253}
K.~Doubleday \textit{et al.}, ``Fuel Cell Inverter Dataset --- GFL, GFM,
and Transition Operation,\'\' NREL Data Repository, DOI: 10.7799/2483511, 2024.

\bibitem{b_unifi}
B.~Kroposki \textit{et al.},
``UNIFI specifications for grid-forming inverter-based resources---Version~1,''
U.S.\ Department of Energy / NREL, Tech.\ Rep., 2022.
[Online]. Available: \url{https://www.energy.gov/sites/default/files/2023-09/Specs\%20for\%20GFM\%20IBRs\%20Version\%201.pdf}

\bibitem{b_ieee1547}
\textit{IEEE Standard for Interconnection and Interoperability of
Distributed Energy Resources with Associated Electric Power Systems
Interfaces}, IEEE Standard 1547-2018, 2018.

\bibitem{b_mcnemar}
Q.~McNemar,
``Note on the sampling error of the difference between correlated
proportions or percentages,''
\textit{Psychometrika}, vol.~12, no.~2, pp.~153--157, 1947.

\bibitem{b_nb06}
L.~G.~C.~Rosa,
``Cross-scenario generalization protocol for ML-based inverter mode
detection,''
\textit{GitHub repository, datacenter-energy-platform, notebook 06},
2024.
[Online]. Available: \url{https://github.com/EngEleLuiz/datacenter-energy-platform}

\end{thebibliography}

\end{document}
% ============================================================================
